\section{Limitations}
\label{sec:gr-limitations}

The benefits of Graph RAG come at the cost of building and maintaining the graph
itself. Constructing a knowledge graph from text requires a multi-stage extraction
pipeline---entity recognition, relation extraction, coreference resolution, and
entity linking---each of which introduces errors that propagate into the
graph~\cite{ji2021kgsurvey, shen2015entitylinking}. Graph-based indexing further
adds up-front cost, since methods that derive an entity graph and pre-generate
community summaries invoke a language model many times over the corpus before any
query is served~\cite{edge2024graphrag}.

The graph must also be kept consistent and current: knowledge graphs are dynamic,
and integrating new or changing data without corrupting existing structure is a
recognized challenge~\cite{hogan2021kg}. Systems such as LightRAG explicitly
introduce incremental update algorithms to address this, which is itself evidence
that maintenance is a first-order concern rather than an afterthought~\cite{guo2024lightrag}.
Finally, ontology design demands domain expertise, and an inadequate schema limits
the consistency and inferential power the graph can provide~\cite{gruber1993ontology,
hogan2021kg}.
